% TOP_PROBLEM: {"id": 3257, "title": "The Bermond-Thomassen Conjecture", "category_id": 3, "category": {"id": 3, "name": "graph_theory", "display_name": "Graph Theory", "description": "Problems involving graphs, networks, and their properties.", "slug": "graph-theory", "order_index": 3, "created_at": "2026-07-31T15:26:25.671Z"}, "status": "open", "problem_number": "OPG-611", "set_id": 12}
% FIRST_POSTED: 2026-10-01
% ENTRY_KIND: statement_only
% UnsolvedMath ID 3257; source catalogue code OPG-611
% !TeX program = lualatex
% Definition and sources only; this file is not an LLM solution attempt.
\documentclass[11pt,a4paper]{article}
\usepackage[margin=25mm]{geometry}
\usepackage{fontspec}
\setmainfont{lmroman10-regular.otf}[BoldFont=lmroman10-bold.otf,
 ItalicFont=lmroman10-italic.otf,BoldItalicFont=lmroman10-bolditalic.otf]
\usepackage{amsmath,amssymb,mathtools}
\usepackage{unicode-math}
\setmathfont{latinmodern-math.otf}
\AtBeginDocument{\let\setminus\smallsetminus}
\usepackage{xurl,hyperref}
\hypersetup{colorlinks=true,urlcolor=blue}
\setcounter{secnumdepth}{0}
\setlength{\parindent}{0pt}
\setlength{\parskip}{5pt}
\setlength{\emergencystretch}{3em}
\begin{document}
\section*{The Bermond-Thomassen Conjecture}\label{3257}
{\small UnsolvedMath ID 3257 · OPG-611 · Graph Theory · OpenGarden}
\subsection{Definitions and mathematical statement}
A finite simple digraph $D=(V,A)$ here has no loops and at most one arc
for each ordered pair of distinct vertices. Both $uv$ and $vu$ are
allowed. The outdegree $d^+(v)$ is the number of arcs with tail $v$, or
equivalently the number of its distinct outneighbours. Set
$\delta^+(D)=\min_{v\in V}d^+(v)$ for a nonempty digraph.

A directed cycle consists of distinct vertices $v_1,\ldots,v_m$ with
$m\ge2$ and arcs $v_iv_{i+1}$, with indices taken modulo $m$. Thus two
opposite arcs form an allowed cycle of length two. Cycles are
vertex-disjoint if their vertex sets have empty pairwise intersection.

The Bermond--Thomassen conjecture is that for every integer $k\ge1$ and
every such digraph,
\[
         \delta^+(D)\ge2k-1
        \quad\Longrightarrow\quad
        D\text{ contains }k\text{ vertex-disjoint directed cycles}.
\]
No indegree condition and no strong-connectivity assumption are imposed.
The cycles may have different lengths and need not cover all vertices.

The threshold cannot be lowered to $2k-2$: the complete symmetric
digraph on $2k-1$ vertices has this outdegree and every directed cycle
uses at least two vertices, precluding $k$ vertex-disjoint cycles.
Here ``symmetric'' means that each pair of vertices has both possible
arcs. This example also explains why confusing vertex disjointness with
arc disjointness would change the sharp question.
\subsection{Short English statement}
Does minimum outdegree at least 2k−1 force k vertex-disjoint directed cycles?
\subsection{Sources}
\begin{itemize}
\item[S1] ULAM.AI and the UnsolvedMath Contributors, supplied problems.json, record 3257 (OPG-611), OpenGarden collection. Catalogue identity and source transcription; CC BY 4.0.
\url{https://huggingface.co/datasets/ulamai/UnsolvedMath}.
\item[S2] Open Problem Garden, The Bermond-Thomassen Conjecture. Original problem and discussion; the mathematical formulation below is expanded from the supplied source record.
\url{http://www.openproblemgarden.org/op/the_bermond_thomassen_conjecture}.
\end{itemize}
\end{document}
